\documentclass[a4paper,11pt]{article}
\usepackage[fontset=none]{ctex}
\setCJKmainfont{Microsoft YaHei}
\setCJKsansfont{Microsoft YaHei}
\setCJKmonofont{Microsoft YaHei}
\usepackage{amsmath,amssymb,graphicx,geometry,booktabs,float,hyperref,listings,xcolor,tabularx,longtable,array,enumitem}
\usepackage[most]{tcolorbox}
\geometry{margin=2.1cm}
\hypersetup{colorlinks=true,linkcolor=blue!55!black,urlcolor=blue!65!black}
\newtcolorbox{knowledgebox}[1]{enhanced,colback=blue!4,colframe=blue!70!black,title=#1,fonttitle=\bfseries,breakable}
\newtcolorbox{importantbox}[1]{enhanced,colback=yellow!8,colframe=orange!80!black,title=#1,fonttitle=\bfseries,breakable}
\newtcolorbox{warningbox}[1]{enhanced,colback=red!4,colframe=red!70!black,title=#1,fonttitle=\bfseries,breakable}
\newtcolorbox{practicebox}[1]{enhanced,colback=green!4,colframe=green!55!black,title=#1,fonttitle=\bfseries,breakable}
\newtcolorbox{boundarybox}[1]{enhanced,colback=purple!4,colframe=purple!65!black,title=#1,fonttitle=\bfseries,breakable}
\lstset{basicstyle=\ttfamily\footnotesize,breaklines=true,frame=single,numbers=left,numberstyle=\tiny,showstringspaces=false,columns=fullflexible}
\setlength{\parindent}{2em}
\setlength{\parskip}{0.38em}
\setcounter{tocdepth}{1}
\renewcommand{\arraystretch}{1.24}
\newcommand{\lecturefigure}[3]{%
  \begin{figure}[H]
  \centering
  \includegraphics[width=0.96\textwidth,height=0.49\textheight,keepaspectratio]{\detokenize{figures/#1}}
  \caption{#2（原视频画面：#3）}
  \end{figure}
}

\begin{document}
\begin{titlepage}
\centering
{\Large 南京大学《生成式软件工程》2026 秋季学期\par}
\vspace{0.8cm}
{\huge\bfseries 第 2 讲：提示词工程\par}
\vspace{0.35cm}
{\Large 从一句话指令到可验证的上下文工程\par}
\vspace{0.55cm}
{\large 主讲：蒋炎岩\qquad 中文精讲笔记\par}
\vspace{0.7cm}
\includegraphics[width=0.90\textwidth,height=0.43\textheight,keepaspectratio]{cover.jpg}\par
\vfill
\begin{tcolorbox}[width=0.93\textwidth,colback=black!2,colframe=black!55]
\textbf{视频标题}：提示词工程 [02-Raw/26生成式软件工程/NJU]\par
\textbf{视频作者/频道}：绿导师原谅你了\par
\textbf{发布日期}：2026-09-01\par
\textbf{时长}：01:38:25\par
\textbf{视频}：\url{https://www.bilibili.com/video/BV1CQt365EzW/}\par
\textbf{课程主页}：\url{https://jyywiki.cn/GSE/2026/}
\end{tcolorbox}
\end{titlepage}

\tableofcontents
\newpage

\section{导读：提示词不是魔法咒语}

\begin{importantbox}{本讲主问题}
当 Agent 可以读取项目、调用工具并持续工作时，我们究竟在“提示”什么？答案不再是一句话，而是一套把意图、事实、权限、行动和验收组织起来的运行环境。本讲真正讨论的是：如何让模型在这个环境里持续向可验证结果推进。
\end{importantbox}

很多提示词教程从句式开始：指定角色、补充背景、要求一步一步思考、给出输出格式。这些方法并非无效，但它们只看见了用户最后输入的那一小段文字。真实 Agent 在第一次回答之前，往往已经读过系统提示词、项目级规则、按需加载的 Skill、工具定义、权限策略、环境状态以及历史对话。用户消息只是其中靠近末端的一部分。

因此，本讲把问题从“怎样写一句更灵的 prompt”推进到“怎样构造一个模型可以可靠行动的 context”。这里的可靠不是语气更像专家，而是四件事同时成立：目标没有歧义，证据保持新鲜，动作边界可执行，完成条件可检查。只追求漂亮输出，会得到一次演示；把这四件事组织成闭环，才可能得到可重复的工程能力。

\begin{knowledgebox}{贯穿全讲的五条线索}
\begin{enumerate}[leftmargin=2.2em]
  \item \textbf{上下文是输入面}：模型实际响应整个 token 序列，而非孤立的一句话。
  \item \textbf{验证器定义终点}：若系统无法识别成功，模型也无法稳定知道何时停止。
  \item \textbf{约束压缩搜索空间}：每个未说明的选择都会由模型按默认偏好补齐。
  \item \textbf{反馈创造长程能力}：尝试、检查、观察和修复组成的循环比一次“深思”更重要。
  \item \textbf{人仍承担外部责任}：人定义问题、边界和后果，AI 在边界内扩大搜索与执行。
\end{enumerate}
\end{knowledgebox}

\subsection{证据与整理方法}

本笔记使用原视频、整段本地中英 ASR、每 15 秒抽取的 394 张完整画面，以及课程主页作为证据。Bilibili 页面未暴露字幕轨，因此字幕是本地转写并做保守术语修正的辅助时间轴，不被当作逐字稿。39 张插图均经过联系表筛选、全分辨率复核和字幕上下文核对；画面保留讲台与投影的原始布局，没有把讲者裁掉后冒充原课件。

\begin{boundarybox}{阅读边界}
讲义中的历史节点、论文案例、漏洞编号和实验数字，首先还原的是课堂陈述。除课程主页外，本笔记没有把这些内容重新包装成独立的外部事实核查。现场生成结果说明的是方法和失败模式，不代表某个模型、仓库或评审流程已经通过通用基准验证。
\end{boundarybox}

\subsection{全讲路线图}

\begin{center}
\small
\begin{tabularx}{\textwidth}{>{\bfseries}p{0.19\textwidth} X}
\toprule
视频区间 & 核心问题 \\
\midrule
00:00--10:30 & Prompt 从哪里来；Agent 看见什么；AI 如何把第一次学习尝试变便宜 \\
10:30--18:45 & 系统提示词、项目规则、Skill、工具、Harness、沙箱和人工批准如何分层 \\
18:45--31:15 & 什么叫完成；如何把 Intent 逐层压缩成 Spec 与可验收的 Impl \\
31:15--45:30 & Token 怎样形成计算深度；形式约束与 Verifier 能看到什么、看不到什么 \\
45:30--54:45 & 长上下文中的注意力竞争；规则该常驻还是按需加载 \\
54:45--71:00 & AI 论文审稿：怎样从角色扮演升级为证据链和独立反证 \\
71:00--84:45 & Skill 蒸馏、Long-Horizon、Greenfield 与参考代码库的帮助/负迁移 \\
84:45--98:25 & 研究失败、吞吐量、机制化记忆、不确定性搜索和多样化采样 \\
\bottomrule
\end{tabularx}
\end{center}

\subsection{本章小结}

提示词工程不是寻找一句万能话术，而是让模型持续看见正确的问题、可用的证据、合法的动作和明确的终点。后续所有例子都可以放回这四格框架中理解。

\section{从 Prompt 到 Context：模型真正看见什么}

\subsection{ChatGPT 普及了术语，但没有发明它}

课堂从“prompt”一词的历史开始。幻灯片把 GPT-2（2019）、GPT-3（2020）以及 2021 年前后的 prompt programming、prompt-based learning 连在一起，说明今天熟悉的提示词概念有连续演化过程。ChatGPT 的关键作用，是把自然语言提示变成大众与模型交互的日常入口，而不是凭空创造了这套思想。

\lecturefigure{fig_001_prompt_origin.jpg}{Prompt 术语早于 ChatGPT：从 GPT-2、GPT-3 到提示词编程}{00:00:15--00:00:30}

聊天框容易制造一个错觉：屏幕上看见的输入框就是模型的全部输入。对 Agent 而言，真正的条件分布更接近幻灯片上的表达：

\[
  \Pr[\text{token}\mid\text{context}]
\]

这里的 context 不仅包含用户文字，也包括系统角色、目标模式、持久规则、已加载的技能、工具调用及返回值。甚至“这段内容是指令、证据还是工具输出”的边界，也由 Harness 以消息类型和优先级提供给模型。

\lecturefigure{fig_002_prompt_beyond_chat.jpg}{提示词不只在对话框，Agent 看到的是整个运行上下文}{00:01:30--00:01:45}

\subsection{上下文工程管理的是工作集}

上下文窗口有容量，却不等于拥有同样大小的有效注意力。把所有文档都塞进去，会同时引入过期事实、冲突指令、无关细节和来源不明的结论。课程给出的上下文工程目标可概括为四点：足够支撑当前动作、随状态更新、能回到来源、能被验证。

\lecturefigure{fig_003_context_engineering.jpg}{从提示词工程转向上下文工程：组织足够、最新、有来源、可验证的工作集}{00:03:00--00:03:15}

\begin{practicebox}{一次上下文盘点}
在让 Agent 开始大任务前，列出六类输入：目标与交付物、权限与安全、工作流与工具、持久规则与风格、环境事实与证据、冲突时的优先级。然后追问：每一项从哪里来？何时更新？谁能覆盖？怎样证明模型真的遵守了？
\end{practicebox}

\begin{warningbox}{“模型说它看见了”不是验证}
让模型复述自己的上下文，只能得到一份候选解释。真正的验证要观察行为：它是否使用了规定语言，是否在权限边界请求批准，是否运行了指定测试，是否保留了不能覆盖的文件。报告是计划，行为才是证据。
\end{warningbox}

\subsection{本章小结}

从 prompt 到 context 的变化，是从句子优化转向系统设计。一个高质量上下文不追求最长，而追求对当前动作最有用、冲突最少、来源最清楚、验证最直接。

\section{立即行动：把 AI 变成学习反馈}

\subsection{第一次尝试变得廉价}

AI 最直接的教育价值，不是替学生交卷，而是降低“先试一下”的成本。过去，一个不确定的问题可能因为环境配置、资料检索和入门代码而被拖延；现在可以立刻构造一个最小例子，让 Agent 解释、实现或运行，再用结果校正自己的理解。课程把这种变化称为行动飞轮：问题越小，反馈越快；反馈越快，下一次提问越具体。

\lecturefigure{fig_004_immediate_action.jpg}{AI 降低第一次尝试成本：从最小问题开始转动反馈飞轮}{00:04:45--00:05:00}

这并不意味着把任务整包外包。若学习者只收集模型答案，他获得的是完成感，不是能力。有效的循环应当保留自己的预测：先说出你认为会发生什么，再让系统运行；观察差异后，解释为什么错；最后改变一个条件重新测试。没有预测的实验很难暴露认知缺口，没有复测的解释也很难形成稳定模型。

\subsection{费曼学习法：让 AI 当听众}

课程把费曼学习法反过来使用：AI 不只是老师，也可以是耐心的听众。学习者选一个概念，用自己的话解释，让 AI 追问不清楚的地方，再回到教材、程序或实验核实。关键动作是“复述”，因为复制原文可以掩盖理解缺口，而面对追问重新组织因果关系会暴露断点。

\lecturefigure{fig_005_feynman_learning.jpg}{让 AI 当听众：复述、暴露缺口、核实、重新讲授}{00:09:30--00:09:45}

\begin{practicebox}{五分钟学习闭环}
\begin{enumerate}[leftmargin=2.2em]
  \item 写下一个足够小、可以在五分钟内验证的问题；
  \item 先给出自己的预测和理由；
  \item 让 AI 提供最小解释或实验，不要一次扩成完整教程；
  \item 用代码、教材或反例核对，标出预测与结果的差异；
  \item 不看原回答，重新讲给 AI，并要求它只追问最薄弱的一步。
\end{enumerate}
\end{practicebox}

\begin{importantbox}{学习证据发生了什么变化}
“答案正确”在 AI 时代越来越弱。更有力的证据是：能否解释关键决策，能否预测变体，能否构造反例，能否在约束变化后修正实现。AI 可以参与每一步，但不能替学习者拥有这些可迁移的判断。
\end{importantbox}

\subsection{本章小结}

AI 让尝试变便宜，但学习仍由反馈质量决定。把模型当作听众、实验搭档和反例生成器，比把它当作答案仓库更能形成长期能力。

\section{上下文的分层：规则、工具与安全边界}

\subsection{先解剖，再优化}

当 Agent 行为异常时，最有效的第一步往往不是再补一句强硬提示，而是盘点它已经看见什么。课程现场让 Agent 把上下文分成目标、权限、工作流、规则、环境事实和优先级，并展示项目级 \texttt{AGENTS.md}、按需加载的 Skill、工具 schema 与会话历史如何逐层进入上下文。

\lecturefigure{fig_006_dissect_context.jpg}{解剖当前上下文：目标、权限、流程、规则、事实与冲突优先级}{00:11:00--00:11:15}

Skill 的意义不只是储存一段更长的 prompt。索引可以常驻，只在任务匹配时加载完整说明；任务结束后，细节可以离开活跃上下文。这样既保留可复用方法，又避免每个任务都为无关说明支付注意力成本。

\subsection{谁判断命令安全}

课程用 Codex 代码和安全案例回答了一个容易混淆的问题：模型可以提出命令，但它不是最终安全边界。Harness 解析动作、匹配批准策略并决定放行、询问或拒绝；操作系统沙箱和权限再提供不可绕过的强制约束。模型判断是有用信号，却不能替代执行层隔离。

\lecturefigure{fig_007_safe_command_layers.jpg}{模型建议、Harness 决策、操作系统强制：安全命令的三层责任}{00:12:30--00:12:45}

现场进一步展示了 Guardian 的分层思路：低风险且语法明确的动作可以走静态路由；边界不清的动作交给专门模型评审；高风险或语义仍不确定的动作请求人工批准。课程提到 CVE-2025-54558，意在说明调用边界和参数解析本身就是安全设计的一部分，而不是用一个漏洞证明所有 Agent 都不安全。

\lecturefigure{fig_008_guardian_layers.jpg}{Guardian 的静态路由、模型评审与人工批准分层}{00:17:15--00:17:30}

\begin{knowledgebox}{Fail-closed 的动作门}
一个可靠动作门至少回答五个问题：动作解析后到底是什么？目标资源是谁？当前主体有什么权限？失败后能否恢复？无法判断时默认发生什么？最后一个问题尤其重要：不确定不应被解释为允许。
\end{knowledgebox}

\begin{warningbox}{不要把礼貌拒绝当作安全}
模型在文本中说“我不会执行危险命令”，并不能建立安全边界。提示词可能被覆盖，语义可能被误判，工具调用也可能把危险藏在参数里。真正的边界应位于模型之外，并保留可审计的批准理由与执行结果。
\end{warningbox}

\subsection{常驻规则与按需能力}

稳定的组织约束适合常驻，例如语言、禁止覆盖用户文件、提交前必须运行测试。任务特有的方法更适合按需加载，例如如何制作课件、怎样审稿、怎样部署某类应用。工具返回值则是短期证据：应在状态变化后更新，而不是作为永远正确的记忆保留。

\begin{center}
\small
\begin{tabularx}{\textwidth}{>{\bfseries}p{0.22\textwidth} p{0.23\textwidth} X}
\toprule
上下文层 & 生命周期 & 适合内容 \\
\midrule
系统/组织规则 & 长期 & 权限边界、沟通风格、不可违反的安全约束 \\
项目规则 & 项目期 & 构建命令、目录边界、代码约定、验收门 \\
Skill & 任务期 & 可复用流程、判断标准、检查点、失败转向 \\
会话证据 & 短期 & 当前文件、日志、工具输出、临时决定 \\
用户动作 & 当前轮 & 目标变化、授权范围、优先级调整 \\
\bottomrule
\end{tabularx}
\end{center}

\subsection{本章小结}

上下文分层同时解决两个问题：让正确规则在需要时出现，让真正的安全边界留在模型之外。写更多提示词不是目标，建立可解释的控制面才是。

\section{什么叫完成：从 Intent 到 Spec 再到 Impl}

\subsection{模型为什么会提前交卷}

语言模型擅长继续一段看起来像答案的文本，却不会天然知道真实系统何时完成。若上下文只出现“继续努力”“做到最好”，模型可能在输出结构看起来完整时结束；若上下文含有可观察的成功谓词，它就能把每次尝试同终点比较。课程把完成条件压缩为：

\[
  V(T)=1
\]

其中 $T$ 是当前产物或轨迹，$V$ 是验收器。公式并不承诺所有价值都能化成布尔值，它强调的是：凡是能明确检查的要求，都不要只留在人的愿望里。

\lecturefigure{fig_009_done_in_context.jpg}{上下文必须描述“完成”是什么，模型才有稳定的停止信号}{00:19:00--00:19:15}

\subsection{把愿望改写成验收条件}

“做一套好用的导航”无法直接判定；“在指定误差内测得经度，并通过独立海试”则提供了证据门槛。课程以 1714 年英国《经度法》说明，工程进步常来自把宏大愿望改写为可验证结果：不是规定每一步怎样做，而是定义什么结果值得奖励。

\lecturefigure{fig_010_verifiable_acceptance.jpg}{把含糊愿望改写为可运行、可观察、可判定的验收条件}{00:21:00--00:21:15}

\begin{practicebox}{把“好”改写成检查口}
先保留意图，再补三类证据：功能证据回答“能不能做”；约束证据回答“是否在权限、性能和格式边界内”；使用证据回答“真实用户能否完成关键任务”。若某项暂时不能自动化，就明确写成需要人工复核，而不是假装它已经可验证。
\end{practicebox}

\subsection{Intent、Spec、Impl 的逐级放大}

一句“生成极品飞车”隐藏了镜头、操控、地图、美术、碰撞、胜负、资产许可和性能目标。Intent 只表达想要什么体验；Spec 决定哪些行为和约束必须成立；Impl 才选择数据结构、引擎、算法与素材。每省略一层，模型就必须自行补齐更多默认值，候选空间近似按各选择维度相乘地扩大。

\lecturefigure{fig_011_intent_spec_impl_gap.jpg}{Intent 到 Spec 再到 Impl：每一次省略都在扩张选择空间}{00:22:30--00:22:45}

\lecturefigure{fig_012_driving_game_demo.jpg}{一句话生成驾驶游戏的现场结果：能运行不等于命中了真实意图}{00:22:45--00:23:00}

课堂给出的 Rule of Thumb 是：写清指令，并承担所有“不明确”的后果。这不是要求写无限长的文档，而是要求选择性地消除高风险歧义。低价值细节可以授权模型决定；会改变产品身份、法律责任、数据安全或验收结果的细节，必须由人明确。

幻灯片以 1854 年 Raglan 的战场命令为反例：发令者默认骑兵与自己拥有相同视野，省略了“那些炮”究竟指什么，模糊处最终由执行者补全。这个例子把“歧义”从语言瑕疵提升为工程风险：执行能力越强，错误默认值造成的后果越大。

\lecturefigure{fig_013_rule_of_thumb.jpg}{要么完成，要么显式失败：不应以道歉或含糊说明代替证据}{00:24:30--00:24:45}

\begin{importantbox}{把失败写进契约}
好的任务契约不仅描述成功，也描述失败：运行哪些检查；检查失败后最多重试几次；缺少证据时如何报告；哪些情况必须停止并请求授权。这样“未完成”成为可接受、可诊断的状态，而不是被模型用漂亮措辞掩盖。
\end{importantbox}

\subsection{本章小结}

Intent 决定方向，Spec 压缩选择空间，Impl 提供可运行产物，Verifier 决定是否到达。提示词工程真正困难的部分，是知道哪些自由度可以授权、哪些必须被规格和证据锁定。

\section{约束不是形容词：从“AI 味”到设计语言}

\subsection{抽象形容词不能排除模板}

“现代、简洁、有科技感”看似具体，实际上几乎没有排除任何模板。模型可以合法地生成渐变标题、玻璃卡片、深色背景，也可以选择完全不同的布局；结果都能自称“现代”。更好的提示不是堆更多风格词，而是指定可组合的设计语言：色彩范围、字体层级、栅格、留白、组件参考、交互状态和禁用模式。

\lecturefigure{fig_014_bad_better_homepage.jpg}{Bad 与 Better：抽象风格词对比可组合、可复核的设计约束}{00:25:30--00:25:45}

\lecturefigure{fig_015_generated_homepage.jpg}{在具体风格、排版与组件约束下生成的个人主页}{00:26:15--00:26:30}

现场主页不是一个通用设计基准，它展示的是“合法解空间”概念。若约束只说个人主页，任何姓名、栏目和视觉系统都可能合规；若给出内容来源、参考组件与截图验收，模型就更容易在可接受空间内搜索。可验证设计仍需要真实浏览器渲染、不同视口、可访问性和内容准确性检查。

\begin{practicebox}{把形容词翻译成设计变量}
把“专业”拆成信息层级、对齐方式、密度和内容证据；把“现代”拆成字体、色彩、间距、动效和组件语言；把“好用”拆成关键任务、状态反馈、错误恢复和可访问性。每个变量都应有参考、范围或截图检查口。
\end{practicebox}

\subsection{人定义问题，AI 搜索实现}

课程提出新的分工：人更靠近 Intent 和 Spec，AI 更擅长在实现空间里搜索。实现成本下降后，真正稀缺的是对产品领域的认识、对用户后果的判断，以及知道什么不能交给默认值。Brooks 在 1986 年讨论软件的 essence 与 accident；课程借此强调，编码偶然性变便宜，并不会让需求、规格和设计概念结构消失。

\lecturefigure{fig_016_human_problem_ai_search.jpg}{人定义问题和验收目标，AI 在实现空间持续搜索}{00:31:15--00:31:30}

\begin{boundarybox}{“AI 搜索实现”不等于人只写需求}
实现过程会暴露新约束，Spec 也会被实验修正。人的工作不是在起点写完一份完美文档后退出，而是维护目标、吸收证据、决定取舍，并对系统进入真实世界后的后果负责。
\end{boundarybox}

\subsection{本章小结}

提示质量的核心不是辞藻，而是选择空间。把审美和产品意图翻译成可组合变量、参考物和验收口，才能让模型的实现能力服务于真实目标。

\section{Token、思维链与 Verifier：浅步骤怎样串成深计算}

\subsection{一步很浅，一串可以很深}

模型生成每个 token 时只做一次有限前向计算，但新 token 会进入后续 context，成为下一步可读取的中间结果。课程把它类比为图灵机纸带：单步转移简单，中间状态被写下后，一串步骤可以积累计算深度。这解释了为何 test-time scaling 有效：不是一句提示突然提升参数，而是给模型更多建立、检查和修正中间状态的机会。

\lecturefigure{fig_017_token_context_depth.jpg}{单步 token 计算有限，但中间结果进入上下文后可以串成深推理}{00:33:15--00:33:30}

多写并不自动等于多想。只有当中间状态能被后续步骤使用，并且存在反馈信号时，额外 token 才可能带来改进。没有检查口的长思维链可能只把最初假设写得更完整；有检查器的短循环则可以不断把错误拉回。

\begin{knowledgebox}{计算深度来自可复用中间状态}
有效中间状态应当让下一步更容易判断：当前假设是什么，已有证据是什么，哪项检查失败，下一次只改变什么。若一段推理没有留下这些结构，它对后续 token 的价值就接近叙述，而非计算。
\end{knowledgebox}

\subsection{形式约束把生成变成搜索}

课程用藏头诗、数字顺序、字数、部首和同音故事展示 Verifier 的力量。开放写作本来有巨大空间；一旦约束可以程序检查，模型就能先生成候选，再根据失败位置继续补偿。它不必一次命中，而可以在上下文中保留“哪一条没过”的局部反馈。

\lecturefigure{fig_018_constrained_poem.jpg}{藏头、数字顺序与八句结构把诗歌生成变成可检查任务}{00:38:45--00:39:00}

\lecturefigure{fig_019_eighteen_char_alliteration.jpg}{十八字、同一偏旁与故事性共同压缩开放写作空间}{00:41:45--00:42:00}

\lecturefigure{fig_020_yi_pinyin_story.jpg}{所有汉字同一读音、仍要讲现代故事的极端约束实验}{00:43:45--00:44:00}

\subsection{通过测试不等于写得好}

Verifier 能看见藏头、字数、顺序、偏旁和拼音，却很难可靠判断趣味、意象、论证力度或文化语感。形式约束可以从搜索空间中排除大量错误候选，却不能覆盖人类价值的全部维度。课程以乔治·佩雷克 1969 年的《消失》说明，严格形式可以催生创造，但作品价值并不等同于“没有使用某个字母”。

\lecturefigure{fig_021_verifier_visible_invisible.jpg}{Verifier 看得见形式约束，却未必看得见趣味、论证与品味}{00:44:30--00:44:45}

\begin{warningbox}{测试套件是代理目标}
模型被训练成执着满足检查时，也可能找到测试的漏洞。测试覆盖表面行为，不能证明用户目标全部满足。需要保留隐藏测试、独立评审、真实使用情境和反例构造，防止“分数被刷满，目标被绕开”。
\end{warningbox}

\begin{practicebox}{双层 Verifier}
第一层使用机器可判定检查，快速排除格式、类型、边界、权限和回归错误；第二层由独立人或独立路线检查解释、审美、证据充分性和真实用户后果。两层结论冲突时，不要用第一层 PASS 覆盖第二层问题。
\end{practicebox}

\subsection{本章小结}

Token 提供连续试错的计算载体，Verifier 提供方向。真正的能力来自两者耦合：中间状态能被复用，错误能被定位，结果能被独立检查，同时承认检查器永远有盲区。

\section{注意力也是资源：规则何时常驻，何时加载}

\subsection{过去一直在场，但不保证一直被用好}

Self-attention 允许每个新 token 重新读取过去上下文。全局的“用中文回答”可能一直影响输出；局部的代码错误也可能在总结阶段重新被注意。然而“可见”不等于“可靠利用”。当模型专注于实现一个子任务时，邻近代码、工具输出和当前错误会占据更强注意力，早期目标和远处约束可能被弱化。

\lecturefigure{fig_022_past_context_attention.jpg}{过去的上下文一直可见，但每次 token 对旧信息的注意并不相同}{00:47:30--00:47:45}

这解释了阶段边界为何重要。实现结束后先做一次“呼吸”：重述目标、记录已证实事实、列出未解决问题，再进入下一阶段。总结不是为了写漂亮报告，而是把关键状态搬到注意力更有利的位置。

\subsection{规则会纠偏，也会抢注意力}

把每次失败都追加进全局规则，会逐渐形成厚重、互相防御的提示墙。规则越多，冲突越难发现，模型在每一步都需要分配注意力处理无关约束。相反，完全不保留经验又会重复踩坑。课程给出的方向是分层：稳定原则常驻，任务方法放 Skill，细节通过测试和接口固化。

\lecturefigure{fig_023_context_competes_attention.jpg}{上下文既是纠偏资源，也是竞争注意力的资源}{00:50:45--00:51:00}

\begin{importantbox}{把“记住”升级为“机制”}
关键约束不要只写在 prompt：类型约束交给编译器，接口约束交给 property tests，回归约束交给 CI，权限边界交给执行层，阶段状态交给 checkpoint。机制可以失败并留下证据；纯文本提醒失败时，模型常常仍会声称自己遵守了。
\end{importantbox}

\begin{center}
\small
\begin{tabularx}{\textwidth}{>{\bfseries}p{0.22\textwidth} X X}
\toprule
问题 & 只追加 prompt & 机制化做法 \\
\midrule
输出格式漂移 & 再强调一次格式 & schema 校验并返回具体字段错误 \\
回归反复出现 & 写“不要破坏旧功能” & 固化回归测试与 CI gate \\
权限越界 & 写“请小心” & 沙箱、allowlist、人工批准 \\
长任务忘记目标 & 重复粘贴完整历史 & checkpoint 保存目标、证据和未决问题 \\
方法难以复用 & 保存上次最终答案 & Skill 保存搜索策略、判断标准和转向条件 \\
\bottomrule
\end{tabularx}
\end{center}

\subsection{本章小结}

上下文窗口是容量，不是保证。通过分层、阶段重述和机制化检查，把最关键的信息放在最能影响下一步的位置，才是长任务中的记忆工程。

\section{案例：把 AI Reviewer 从角色扮演升级为证据审计}

\subsection{“你是审稿人”只规定了表面角色}

课程现场把自己的论文交给 AI，要求模拟 ICSE 审稿。简单角色提示容易学到审稿格式、语气和常见栏目，却不保证模型执行了真正的证据程序。它可能复述摘要、罗列一般性优缺点，再给出看似合理的分数；这种输出像 review，却没有建立 claim 到 evidence 的链条。

\lecturefigure{fig_024_ai_peer_review.jpg}{让 AI 审自己的论文：角色提示不等于掌握审稿方法}{00:55:00--00:55:15}

\lecturefigure{fig_025_paper_review_demo.jpg}{现场读取论文并生成审稿意见：可观察轨迹比最终语气更重要}{00:55:45--00:56:00}

更可靠的做法是先提供审稿方法：论文声称的主要贡献是什么；每个核心 claim 需要什么证据；实验究竟排除了哪些替代解释；哪些结论超出数据覆盖；若站在反方，最小反例是什么。身份可以帮助设定受众和格式，方法才决定推理路径。

\subsection{显式建立 claim--evidence--assumption}

现场演示把论文断言、讲者记忆和置信度放入表格。表格的价值不是让 AI “相信作者”，而是暴露证据层级：哪些事实来自论文，哪些来自已有知识，哪些只是需要进一步查证的记忆。对于论文外部事实，应回到原始来源；对于实验结论，应追到表格、脚本和数据，而不是停在摘要措辞。

\lecturefigure{fig_026_memory_claim_table.jpg}{把论文断言、已有记忆与置信度对齐，显式暴露证据缺口}{01:00:30--01:00:45}

\begin{practicebox}{一条 claim 的审计卡}
\begin{enumerate}[leftmargin=2.2em]
  \item 原文到底声称什么，范围与条件是什么；
  \item 需要哪类证据才能支持；
  \item 论文实际提供了什么，证据位于何处；
  \item 哪些替代解释没有排除；
  \item 什么观察会推翻或显著削弱结论；
  \item 当前结论的置信度以及仍需补查的来源。
\end{enumerate}
\end{practicebox}

\subsection{三条独立路线先分离，再汇合}

课程建议把审稿拆成三条相对独立的路线。第一条剥离贡献：先不看作者包装，重写问题、方法和新增信息；第二条证据审计：为每个核心 claim 建账本；第三条反方复现：在最弱或最反直觉的位置重算结果、寻找最小反例。三条路线独立形成判断后再汇合，可以减少多个 reviewer 只是重复同一锚点的风险。

\lecturefigure{fig_027_three_review_routes.jpg}{剥离贡献、证据审计、反方复现三条独立审稿路线}{01:09:45--01:10:00}

\subsection{不要被作者的方法论带着走}

论文信息密度很高，顺序本身会分配注意力。若 reviewer 从标题、摘要、相关工作一路顺读，作者的术语和因果结构会成为默认框架。课程提醒，critical thinking 可能在不知不觉中变成“在作者框架里挑小问题”。更强的做法是在深读前独立写下：它声称解决什么、最强替代解释是什么、什么结果会推翻结论；再回到正文寻找证据。

\lecturefigure{fig_028_method_limits.jpg}{避免被作者结构锚定：先独立重建 claim、evidence 与 assumption}{01:11:00--01:11:15}

\begin{warningbox}{AI 审稿的可靠边界}
模型可以扩大检索、整理和反例生成，但不能凭流畅语气补全缺失实验。若没有论文、补充材料、代码和数据，它只能评审可见证据。若外部事实没有原始来源，必须标为待核实，而不是利用模型记忆给出确定结论。
\end{warningbox}

\subsection{本章小结}

高质量 Reviewer 不是一种角色语气，而是一套可审计程序。把贡献、证据、假设和反例分离，把不确定性写出来，AI 才能成为研究判断的放大器而不是结论生成器。

\section{从一次成功到 Long-Horizon：Skill、反馈与代码库}

\subsection{把品味蒸馏成 Skill}

当一次任务成功时，最值得保存的不是最终答案，而是通向答案的路线：如何拆解，先看什么证据，用什么判断标准，在哪些检查点停下来，失败后怎样换路。课程称之为把“品味”蒸馏成 Skill。Skill 本身也应形成闭环，避免长指令中后半段越来越难遵循。

\lecturefigure{fig_029_skill_distillation.jpg}{从偶然成功到可重复方法：保存搜索策略、判断标准与转向条件}{01:15:00--01:15:15}

\begin{knowledgebox}{Skill 应保存什么}
保存问题分解、证据路由、判断门、验证方法和失败转向；少保存某次任务的固定答案。把长 Skill 分成常驻目标、任务阶段和按需参考，让每一阶段都有可观察的完成条件。
\end{knowledgebox}

\subsection{长程能力是反馈循环训练出来的}

Long-Horizon 任务包含未知：起点只有大致 Intent，Spec 和 Impl 都会在探索中变化。模型之所以能持续写代码，并非一次规划覆盖了全部未来，而是代码世界提供密集反馈：构建、测试、benchmark、运行日志和界面截图不断告诉系统下一步是否更接近目标。

\lecturefigure{fig_030_long_horizon_feedback.jpg}{尝试、编译/测试、观察、修复：反馈循环形成长程能力}{01:15:45--01:16:00}

反馈也可能教坏模型。若测试只覆盖表面行为，Agent 会把“坚持”变成 reward hacking：反复修改直到分数通过，却绕开真实目标。因此还要保留 held-out tests、独立 review 和不由执行者控制的验收证据。

\subsection{Greenfield 为什么未必更容易}

从零开始的代码很少，看似上下文干净；但所有命名、分层、错误处理、测试和演化边界都尚未决定，选择空间反而最大。成熟仓库含有大量代码，也含有可供 in-context learning 的工程惯例。问题不在字节多少，而在信息密度：哪些代码表达稳定设计，哪些只是历史事故。

\lecturefigure{fig_031_greenfield_paradox.jpg}{Greenfield 代码少不等于信息少；成熟仓库提供惯例也积累噪声}{01:17:00--01:17:15}

课程讨论了“给 AI 一个相近优质项目是否更可靠”的猜想，并引用 RepoCoder、CodeRAG-Bench 等研究线索。幻灯片提到 RepoCoder 的迭代检索相对当前文件基线提升超过 10\%，同时明确保留边界：这支持“好示例可能有用”，却没有证明外来仓库对所有项目都有益；外部有效性、负迁移和检索到错误片段的问题仍未解决。因此，更有价值的是设计对照实验，而不是把“参考仓库越大越好”当成定律。

\lecturefigure{fig_032_reuse_repo_experiment.jpg}{用固定任务、模型和预算检验参考仓库的帮助与负迁移}{01:21:00--01:21:15}

\begin{practicebox}{一个可执行的仓库复用对照}
固定任务、模型、token 预算和随机种子，比较三组上下文：无示例；完整参考仓库；人工精选的 3--5 条 golden paths。记录测试通过率、首次架构违规、返工轮次、review 时间和后续缺陷。只有这样，才能区分“内容更多”与“真正提供了高质量惯例”。
\end{practicebox}

\subsection{本章小结}

Long-Horizon 的核心不是让模型连续运行更久，而是让每一阶段都产生可归因反馈。Skill 保存方法，代码库提供惯例，测试与独立评审防止反馈被投机利用。

\section{研究失败：Token 花完，知识为什么没有增加}

\subsection{总分不是可用反馈}

课程讲述一个符号执行引擎优化案例：Agent 可以不断尝试 benchmark、修改实现、再跑一次，花费大量 token 却没有形成新知识。问题不一定是模型不够聪明，而是反馈太昂贵、太局部不可见、太难归因。一个总 benchmark 分数几乎不告诉系统：慢在路径爆炸、求解器、缓存、内存、搜索策略，还是实现 bug 与环境噪声。

\lecturefigure{fig_033_research_failure.jpg}{Research Failure：token 花完但知识没有增加，原因是缺少局部可归因反馈}{01:25:15--01:25:30}

研究允许“原始想法不成立”成为最佳结论，但 Agent 往往被训练成必须完成任务。没有可证伪假设时，它可能把错误路线包装得越来越完整。课程提到 KLEE 对搜索策略、约束缓存和独立约束集等机制的拆分，强调应先把大问题分成可测机制，再逐项建立几十秒内返回的 micro-benchmark。

\begin{importantbox}{Research Failure 的正确产物}
一次失败至少应增加一种知识：哪类输入触发问题，哪个机制被排除，哪个指标与目标无关，什么实验能区分两个假设。若只留下“跑了很多 token，没有提升”，过程就没有形成可复用证据。
\end{importantbox}

\subsection{少干预，反而做得更多}

人盯着每一步，会优化单个任务的 latency：看到偏航立即纠正，当前任务可能更快。但全天生产力关心 throughput：多个任务能否并行，失败能否恢复，结果能否稍后 review。只要任务边界清楚、失败可恢复，放手让 Agent 跑，再按检查口批量审阅，往往比持续干预更高效。

\lecturefigure{fig_034_less_intervention.jpg}{区分单任务延迟与全天吞吐量：少干预的前提是并行和可恢复}{01:28:45--01:29:00}

“少干预”不是不 review。Review 的重点应从观察每个 token，转向检查接口是否稳定、权限是否越界、数据格式是否正确、测试是否可信、失败能否回滚。丰田式自动化的启发是让异常触发停止和告警，而不是让人持续盯住每一步。

\begin{practicebox}{放手运行前的四个门}
任务是否可以与其他工作并行？错误是否能通过版本控制或快照恢复？执行权限是否被限制在明确范围？结束时是否有足够证据进行独立 review？四项任一为否，持续人工观察仍可能是必要控制。
\end{practicebox}

\subsection{把要求从记忆升级为机制}

随着上下文增长，指令遵循会下降。最后的 prompt 和新加载的 Skill 常更容易获得注意，但这不是可靠记忆方案。课程建议把关键要求重新落地：短 checklist 放在阶段边界，Spec 与 Skill 保留目标和转向，类型、属性测试、CI gate 与权限沙箱提供机器约束，压缩摘要只保留目标、证据和未决问题。

\lecturefigure{fig_035_context_memory_mechanism.jpg}{Context 越长不等于记得越牢：把关键要求从文字升级成机制}{01:31:30--01:31:45}

\subsection{本章小结}

Token 只有在反馈能减少不确定性时才转化为知识。研究任务需要便宜、局部、可归因的实验；生产任务需要并行、可恢复的执行和独立验收；长上下文则需要机制化记忆。

\section{在不确定性中搜索：算力、宽度与多样性}

\subsection{这次不是解题}

封闭问题给定要求、路线和正确答案，可以沿一条思维链逐步逼近。开放研究没有唯一解，也没有预先知道的“已经够好”概率。课程用 Herbert Simon 的“满意化”提醒：搜索必须有停止条件，但停止不是证明找到了全局最优，而是在资源、证据和风险下接受当前方案。

\lecturefigure{fig_036_uncertainty_search.jpg}{开放任务不是解题：显式定义探索路线、停止条件与校验宽度}{01:33:15--01:33:30}

单条 chain-of-thought 用算力换深度；并行生成多条路线，再用独立校验筛选，则用算力换宽度。对于命名、设计、研究假设等任务，只问“给我几个名字”会得到模型偏好下的窄样本。更好的做法是先定义探索轴，再让不同 Agent 或不同种子覆盖组合空间，最后使用另一条路线检查重复、缺陷和隐含偏见。

\lecturefigure{fig_037_compute_for_intelligence.jpg}{把灵感生成改造成并行搜索：先定义空间，再采样和校验}{01:34:30--01:34:45}

\begin{knowledgebox}{深度、宽度与校验的分工}
深度用于沿一条路线解决依赖问题；宽度用于覆盖互斥假设和风格；校验用于淘汰看似新颖但不满足条件的候选。只增加深度会被早期假设锁定，只增加宽度会产生大量噪声，没有独立校验则无法把算力转成判断。
\end{knowledgebox}

\subsection{多采样不是多数投票}

人的未来来自教育背景、生活经验和偏见的多样性。模型多采样若共享同一提示、同一资料和同一评估者，可能只是同一偏见的细微抖动。课程借 Richard Sutton 的 Big World Hypothesis 强调，真实世界远大于单个模型的已知结构；群体智能不要求每个成员无偏，而要求系统不只保留一种偏见，并允许互相质疑。

\lecturefigure{fig_038_diverse_sampling.jpg}{通过不同生成路线与交叉批判保持群体多样性}{01:36:45--01:37:00}

\begin{practicebox}{把灵感任务改造成搜索任务}
先列探索维度和必须满足的约束；让不同采样分别覆盖不同维度，不许先互看答案；用独立校验器查重、检查约束和寻找反例；最后由人比较无法自动量化的意义、风险与长期后果。校验者与生成者若完全共享上下文，应额外引入外部证据或人工复核。
\end{practicebox}

\subsection{墙内的超人，墙外的世界}

结尾的“墙内超人”图像总结了课程立场：在规格和验证器围出的墙内，算力可以把搜索、实现和修复放大到超人的规模；墙外仍是开放世界，包含未写入上下文的人、制度、伦理、资源和不可逆后果。课程引用 Stanford AI Index 2025 中产业来源模型接近九成的统计，提出算力正在成为生产资料，但算力的所有者并不会自动获得正确目标。

\lecturefigure{fig_039_wall_superhuman.jpg}{算力成为生产资料后，Agent 可以在墙内超人化，人仍负责墙和外部世界}{01:37:30--01:37:45}

\begin{importantbox}{最终原则}
能验证的，尽量写成机制；不能验证的，明确保留人的判断；有风险的，把边界放在模型之外；有不确定性的，用多路线搜索而不是伪装成确定答案。提示词工程的成熟标志，不是模型更听话，而是系统更诚实地知道自己完成了什么、没有完成什么。
\end{importantbox}

\subsection{本章小结}

算力可以同时换取推理深度和探索宽度，但只有独立校验、真实证据与多样化路线才能把算力转成可靠判断。人最终负责定义搜索空间、停止条件和不可跨越的墙。

\section{结语：把一句话升级为可验证系统}

本讲从一个日常词“提示词”出发，最终抵达生成式软件工程的控制面。模型看到的是整个 context；context 里的约束决定搜索空间；Verifier 决定系统是否知道自己到达；Harness 和操作系统决定动作能否越界；反馈循环决定长程工作是否真正增加知识。

\begin{knowledgebox}{一份可直接使用的 Agent 任务骨架}
\begin{enumerate}[leftmargin=2.2em]
  \item \textbf{Intent}：为什么做，谁会使用，错误会伤害谁；
  \item \textbf{Spec}：必须成立的行为、不可违反的边界、允许模型决定的自由度；
  \item \textbf{Evidence}：事实来自哪里，何时获取，哪些仍是待核实假设；
  \item \textbf{Actions}：允许使用哪些工具，哪些动作需批准，怎样回滚；
  \item \textbf{Verifier}：机器检查、独立 review、真实使用证据分别是什么；
  \item \textbf{Stop}：成功、显式失败和请求人工介入各自满足什么条件；
  \item \textbf{Memory}：哪些经验沉淀进 Skill，哪些约束升级为测试、接口和 CI。
\end{enumerate}
\end{knowledgebox}

这套骨架也说明，人类工作没有被一个更长 prompt 消灭。真正困难的是识别未写出的目标、选择值得搜索的问题、设计不易被投机的证据、控制真实世界权限，并在证据不足时诚实地停下来。Agent 可以把墙内的执行能力推到极高，但墙如何画、门何时开、墙外的人承担什么后果，仍然是软件工程的核心。

\begin{practicebox}{课后练习}
选一个你最近交给 AI 的真实任务，不增加形容词，只重写它的 Intent、Spec、Evidence、Actions、Verifier、Stop 和 Memory。让另一个人只凭这七项判断：成功是否可识别，失败是否可恢复，权限是否过宽，哪些价值仍需人工判断。再让 Agent 执行最小版本，并记录第一次偏航来自哪一层。
\end{practicebox}

\begin{boundarybox}{资料状态说明}
制作本讲义时，课程主页可确认课程名称、教师与整体安排，但未提供本讲可独立下载的官方课件或逐字字幕。因此，画面内容以原视频为准，字幕仅用于定位与语义互证；所有可编辑源文件、图像清单和质量报告随讲义一并保存，便于未来用官方材料更新。
\end{boundarybox}

\end{document}
